% =============================================================
% 01_enterprise.tex — 企業活動・経営戦略
% =============================================================
\chapter{企業活動と経営戦略}

\chaptermeta{%
  \item 経営理念→方針→戦略→戦術の4段階を説明できる
  \item CxO の各役職の違いを即答できる
  \item SWOT/VRIO/PPM 等のフレームワークを使い分けられる
  \item DX/PoC/リーンスタートアップの関係を整理できる
  \item 新事業・起業の用語を正しく識別できる
}{75分}{2}{なし（本書の出発点）}

Iパス最大の出題分野が「ストラテジ系」です。
本章は\textbf{企業活動 → 経営戦略フレームワーク → IT活用}の順に進みます。
すべて\textbf{「経営者の意思決定」}という1本の軸で貫かれています。

% =============================================================
\section{企業活動と組織}
% =============================================================

\begin{pitfall}[「経営理念」と「経営戦略」がごちゃごちゃ]
抽象から具体への4段階を固定しましょう。
\begin{itemize}
  \item \termtag{経営理念}：企業が存在する根本の目的（なぜ）
  \item \termtag{経営方針}：理念を実現する大きな方向性（どちらへ）
  \item \termtag{経営戦略}：方針を実現する具体的な道筋（どうやって）
  \item \termtag{戦術}：戦略を実現する現場の施策（いま何を）
\end{itemize}
\end{pitfall}

\subsection{組織形態の整理}

\begin{hikaku}[代表的な組織形態]
\comptable{形態 & 特徴 & 強み／弱み}{
職能別組織 & 機能（営業・製造等）で分ける & 専門性↑／全体最適↓ \\
事業部制組織 & 事業ごとに独立採算 & 機動力↑／重複↑ \\
マトリックス組織 & 機能×プロジェクトの2軸 & 柔軟性↑／指揮系統が二重 \\
カンパニー制 & 事業を社内分社化 & 責任明確／小回り困難
}
\end{hikaku}

\subsection{経営層の役職}

\begin{teigi}[CxO：頭文字で役割がわかる]
\termtag{CEO}（最高経営責任者）／\termtag{COO}（執行）／\termtag{CFO}（財務）／
\termtag{CIO}（情報）／\termtag{CTO}（技術）／\termtag{CISO}（情報セキュリティ）。
「C = Chief、O = Officer、真ん中が担当分野」と覚える。
\end{teigi}

% --- 例題 ---
\begin{reidai}{CxOの識別}
技術戦略の策定や技術開発の推進といった技術経営に直接の責任をもつ役職はどれか。

\sentakushi{CEO}{CFO}{COO}{CTO}
\end{reidai}

\begin{shishin}
「技術経営」「技術開発の推進」というキーワードに注目する。
各CxOの真ん中のアルファベットが担当分野を示すので、「T = Technology」を探す。
\end{shishin}

\begin{kaitou}
\textbf{正解：エ（CTO）}

CTO = Chief \textbf{Technology} Officer（最高技術責任者）。
CEO は経営全体、CFO は財務、COO は業務執行の責任者。
\end{kaitou}

\begin{minuki}[見抜き方]
\begin{itemize}
  \item CxOは「真ん中のアルファベット」が守備範囲を示す
  \item 問題文のキーワード（技術/財務/情報/セキュリティ）と照合するだけ
  \item CIO（情報）とCTO（技術）の混同に注意——CIOは「情報\textbf{システム戦略}」、CTOは「\textbf{技術}開発」
\end{itemize}
\end{minuki}

% =============================================================
\section{経営戦略のフレームワーク}
% =============================================================

\begin{pitfall}[横文字フレームワークが多すぎる]
「\textbf{何を見るための道具か}」を決めて引き出しに収めれば怖くありません。
\begin{itemize}
  \item 自社を見る → SWOT、VRIO
  \item 市場を見る → 3C、5F（5フォース）
  \item 売り方を設計 → 4P／4C
  \item 投資配分を決める → PPM
  \item 成長方向を選ぶ → アンゾフの成長マトリクス
\end{itemize}
\end{pitfall}

\begin{hikaku}[戦略フレームワークの引き出し]
\comptable{フレーム & 観点 & 一言で}{
SWOT & 自社の健康診断 & 強み・弱み×機会・脅威 \\
VRIO & 強みの深堀り & 価値・希少性・模倣困難性・組織 \\
3C & 市場の三角測量 & 顧客・競合・自社 \\
5F & 業界の圧力 & 既存・新規・代替・売手・買手 \\
4P/4C & マーケティング & 売り手4視点／買い手4視点 \\
PPM & 投資配分 & 成長率×占有率の4象限 \\
アンゾフ & 成長方向 & 既存/新規 × 市場/製品
}
\end{hikaku}

\subsection{PPMの4象限}

\begin{teigi}[PPM（プロダクト・ポートフォリオ・マネジメント）]
\textbf{市場成長率}（縦軸）と\textbf{市場占有率}（横軸）で事業を4象限に分類する。
\begin{itemize}
  \item \termtag{花形（Star）}：成長率高・占有率高 → 投資継続
  \item \termtag{金のなる木（Cash Cow）}：成長率低・占有率高 → 収穫
  \item \termtag{問題児（Question Mark）}：成長率高・占有率低 → 選別投資
  \item \termtag{負け犬（Dog）}：成長率低・占有率低 → 撤退検討
\end{itemize}
\end{teigi}

% --- 例題 ---
\begin{reidai}{VRIO分析の識別}
投資の優先度などの経営の戦略を策定するために、経済価値、希少性、模倣困難性及び組織の四つの要素で評価することによって、自社のもつ資源を分析する手法として、最も適切なものはどれか。

\sentakushi{4P}{PPM}{SWOT分析}{VRIO分析}
\end{reidai}

\begin{shishin}
問題文に「経済価値、希少性、模倣困難性、組織」と4要素が列挙されている。
各フレームワークの「何を見る道具か」を思い出し、この4要素に一致するものを選ぶ。
\end{shishin}

\begin{kaitou}
\textbf{正解：エ（VRIO分析）}

VRIO = \textbf{V}alue（経済価値）・\textbf{R}areness（希少性）・\textbf{I}nimitability（模倣困難性）・\textbf{O}rganization（組織）。
4P はマーケティングミックス、PPM は成長率×占有率、SWOT は強み弱み×機会脅威。
\end{kaitou}

\begin{minuki}[見抜き方]
\begin{itemize}
  \item 問題文に「4つの要素」が列挙されていたら、各フレームワークの構成要素と照合する
  \item VRIO は「その強みは\textbf{本当に競争優位か}」を4つの問いで確認するツール
  \item SWOT は「内部×外部」の2軸、VRIO は「強みの質」を掘り下げる点が違う
\end{itemize}
\end{minuki}

% =============================================================
\section{新事業とIT活用}
% =============================================================

\begin{hikaku}[新規事業の関連用語]
\comptable{用語 & 要点 & 混同しやすい用語}{
PoC & 技術の実現性確認 & PoV（価値の実証）より前段階 \\
PoV & ビジネス価値の実証 & PoCの次のステップ \\
MVP & 検証用の最小限の製品 & 本番製品ではない \\
リーンスタートアップ & MVPで仮説検証を繰り返す手法 & DX（概念）とは階層が違う \\
DX & ITで経営を変革する概念 & 手法ではなく変革そのもの \\
ハッカソン & 短期集中でプロトタイプを作り競う & セキュリティキャンプとは別
}
\end{hikaku}

% --- 例題 ---
\begin{reidai}{PoCの識別}
新しい概念やアイディアの実証を目的とした、開発の前段階における検証を表す用語はどれか。

\sentakushi{CRM}{PoC}{RAS}{SLA}
\end{reidai}

\begin{shishin}
「概念の実証」「開発の前段階」というキーワードに注目。
Proof of Concept = 概念実証。CRMは顧客管理、RASは信頼性指標、SLAはサービス品質合意。
\end{shishin}

\begin{kaitou}
\textbf{正解：イ（PoC）}

PoC（Proof of Concept）は「概念実証」。新技術やアイデアが実現可能かを開発前に小規模に検証する活動。
\end{kaitou}

% =============================================================
\section{過去問ドリル}
% =============================================================

\shoumatsuhead{過去問ドリル：ストラテジ系}

\begin{itpmon}[\sourcebadge{R7}\enspace\diffbadge{標}]{%
ハッカソンに関する記述として、最も適切なものはどれか。}
\sentakushi
  {定められたルールの下、ある主題について、肯定派と否定派といった異なる立場に分かれて議論する。}
  {情報セキュリティ分野で活躍したいという意志をもった若者が、合宿形式で実践的な知識を学ぶ。}
  {プログラマーやデザイナーなどから成る複数の参加チームが、与えられたテーマに関するプロトタイプを短期間で作成し、その成果を発表して競い合う。}
  {問題解決や利用者獲得などゲーム的な要素のない分野に、デジタル技術を活用したゲームの要素を取り入れることによって、利用者の参加を動機づける。}
\end{itpmon}
\begin{kaisetsu}{ウ}
\textbf{ハッカソン} = hack + marathon。短期集中でプロトタイプを作り競い合うイベント。
アはディベート、イはセキュリティキャンプ、エはゲーミフィケーションの説明。
\end{kaisetsu}

\begin{itpmon}[\sourcebadge{R6}\enspace\diffbadge{標}]{%
未来のある時点に目標を設定し、そこを起点に現在を振り返り、目標実現のために現在すべきことを考える方法を表す用語として、最も適切なものはどれか。}
\sentakushi
  {PoC（Proof of Concept）}
  {PoV（Proof of Value）}
  {バックキャスティング}
  {フォアキャスティング}
\end{itpmon}
\begin{kaisetsu}{ウ}
\textbf{バックキャスティング}は未来の目標から逆算して現在の行動を決める手法。SDGsなどの長期目標設定で使われる。フォアキャスティングは現状から積み上げる方法で、方向が逆。
\end{kaisetsu}

\begin{itpmon}[\sourcebadge{R6}\enspace\diffbadge{標}]{%
ベンチャーキャピタルに関する記述として、最も適切なものはどれか。}
\sentakushi
  {新しい技術の獲得や、規模の経済性の追求などを目的に、他の企業と共同出資会社を設立する手法}
  {株式売却による利益獲得などを目的に、新しい製品やサービスを武器に市場に参入しようとする企業に対して出資などを行う企業}
  {新サービスや技術革新などの創出を目的に、国や学術機関、他の企業など外部の組織と共創関係を結び、積極的に技術や資源を交換し、自社に取り込む手法}
  {特定された課題の解決を目的に、一定の期間を定めて企業内に立ち上げられ、構成員を関連部門から招集し、目的が達成された時点で解散する組織}
\end{itpmon}
\begin{kaisetsu}{イ}
\textbf{ベンチャーキャピタル（VC）}は成長見込みのあるスタートアップに出資し、株式売却（IPO等）で利益を得る投資会社。
アはジョイントベンチャー、ウはオープンイノベーション、エはタスクフォースの説明。
\end{kaisetsu}

\begin{itpmon}[\sourcebadge{R6}\enspace\diffbadge{強}]{%
企業の戦略立案やマーケティングなどで使用されるフェルミ推定に関する記述として、最も適切なものはどれか。}
\sentakushi
  {正確に算出することが極めて難しい数量に対して、把握している情報と論理的な思考プロセスによって概数を求める手法である。}
  {特定の集団と活動を共にしたり、人々の動きを観察したりすることによって、慣習や嗜好、地域や組織を取り巻く文化を類推する手法である。}
  {入力データと出力データから、その因果関係を統計的に推定する手法である。}
  {有識者のグループに繰り返し同一のアンケート調査とその結果のフィードバックを行うことによって、ある分野の将来予測に関する総意を得る手法である。}
\end{itpmon}
\begin{kaisetsu}{ア}
\textbf{フェルミ推定}は、正確な数値を得ることが困難な量を論理的な分解と概算で推定する手法。
イはエスノグラフィ、ウは回帰分析、エはデルファイ法の説明。
\end{kaisetsu}

% =============================================================
\section{タイムアタック}
% =============================================================

\begin{timeattack}{ストラテジ用語クイックチェック}{5分}{6問中5問}

\begin{itpmon}[\diffbadge{基}]{PPMで「市場成長率が低く、市場占有率が高い」事業を何と呼ぶか。}
\sentakushi{花形}{金のなる木}{問題児}{負け犬}
\end{itpmon}
\begin{kaisetsu}{イ}
成長率低・占有率高 = \textbf{金のなる木}。投資を抑えて利益を回収するフェーズ。
\end{kaisetsu}

\begin{itpmon}[\diffbadge{基}]{SWOT分析で扱うのはどれか。}
\sentakushi
  {市場の成長率と占有率}
  {自社の強み・弱みと、外部の機会・脅威}
  {製品・価格・流通・販促の4要素}
  {顧客・競合・自社の3視点}
\end{itpmon}
\begin{kaisetsu}{イ}
アはPPM、ウは4P、エは3C。SWOTは「内部（強み弱み）×外部（機会脅威）」の4象限。
\end{kaisetsu}

\begin{itpmon}[\diffbadge{基}]{DXの説明として最も適切なものはどれか。}
\sentakushi
  {MVPで仮説検証を繰り返す手法}
  {ITを使って経営やビジネスモデルを変革すること}
  {短期集中でプロトタイプを作り競う}
  {概念や技術の実現可能性を検証すること}
\end{itpmon}
\begin{kaisetsu}{イ}
DX = Digital Transformation。アはリーンスタートアップ、ウはハッカソン、エはPoC。
\end{kaisetsu}

\begin{itpmon}[\diffbadge{基}]{PoCとPoVの正しい関係はどれか。}
\sentakushi
  {PoCは価値の実証、PoVは概念の実証}
  {PoCは概念の実証、PoVは価値の実証}
  {どちらも同じ意味}
  {PoCは市場調査、PoVは技術検証}
\end{itpmon}
\begin{kaisetsu}{イ}
PoC = Proof of \textbf{Concept}（概念実証）。PoV = Proof of \textbf{Value}（価値実証）。PoCが先、PoVが後。
\end{kaisetsu}

\begin{itpmon}[\diffbadge{標}]{フォーラム標準の説明として適切なものはどれか。}
\sentakushi
  {公的機関が策定した規格}
  {市場で事実上の標準となったもの}
  {特定分野の企業が集まって策定した規格}
  {個人が提案して広まったもの}
\end{itpmon}
\begin{kaisetsu}{ウ}
フォーラム標準 = 業界の企業団体が合意して作る規格。アはデジュール標準、イはデファクト標準。
\end{kaisetsu}

\begin{itpmon}[\sourcebadge{R6}\enspace\diffbadge{標}]{%
技術戦略の策定や技術開発の推進といった技術経営に直接の責任をもつ役職はどれか。}
\sentakushi{CEO}{CFO}{COO}{CTO}
\end{itpmon}
\begin{kaisetsu}{エ}
\textbf{CTO}（Chief Technology Officer）= 最高技術責任者。技術経営に直接の責任を持つ。
\end{kaisetsu}

\end{timeattack}

% =============================================================
\section{よくあるミスと到達確認}
% =============================================================

\begin{yokumiss}[この章でよくあるミス]
\begin{enumerate}[leftmargin=1.6em, itemsep=5pt]
  \item \textbf{SWOTとVRIOを混同する}——SWOTは4象限マトリクス、VRIOは強みの深堀り。目的が違う。
  \item \textbf{PoCとMVPを同一視する}——PoCは「技術の実現可能性確認」、MVPは「最小限の製品で市場仮説を検証」。
  \item \textbf{DXとリーンスタートアップを混同する}——DXは概念、リーンスタートアップは手法。
  \item \textbf{CIOとCTOを逆にする}——CIOは「情報システム戦略」、CTOは「技術開発」。
\end{enumerate}
\end{yokumiss}

\begin{tatsaku}
  \item 経営理念→方針→戦略→戦術の4段階を説明できる
  \item CxOの各役職の守備範囲を即答できる
  \item SWOT/VRIO/3C/5F/4P/PPMの「何を見る道具か」を言える
  \item PPMの4象限を図なしで再現できる
  \item PoC/PoV/MVP/リーンスタートアップの違いを説明できる
  \item バックキャスティングとフォアキャスティングの違いを言える
  \item ハッカソン/ゲーミフィケーション/セキュリティキャンプを区別できる
  \item ベンチャーキャピタルとジョイントベンチャーの違いを説明できる
\end{tatsaku}

\nextchapter{%
会計・財務・法務を学びます。PLの5段階利益、損益分岐点、知的財産権、契約形態——
ストラテジ系の後半を一気に攻略します。
}
